% !TEX root = ../main.tex
\chapter{Conclusions i treball futur}\label{ch:conclusions}
\fancyhead[RE]{\slshape \textsf{Capítol \thechapter.\ Conclusions}}
\fancyhead[LO]{\slshape \textsf{\nouppercase \rightmark}}

\section{Resposta al problema plantejat}\label{sec:final_answer}

Aquest treball partia d'un problema operatiu: un satèl·lit no pot dependre
d'un operador que dibuixi una ROI abans de cada compressió. Alhora, aplicar
qualitat uniforme a tota la imatge desaprofita bits quan només una part és
prioritària. La solució desenvolupada combina tres elements separables:

\begin{enumerate}[leftmargin=*]
  \item un preset C identifica automàticament blocs d'interès;
	  \item una política C assigna qualitat a ROI i fons;
	  \item el compressor i el descompressor C permeten mesurar error, aplicar el
	  Q-map resultant i reconstruir la imatge sense conèixer la semàntica que l'ha
	  generat.
\end{enumerate}

La separació és la contribució de disseny principal. Permet canviar l'objectiu
de missió sense modificar el còdec: vegetació, contrast, aigua o núvols són
detectors; \code{focus} i \code{preserve-roi} són regles de qualitat. El Q-map
és el contracte entre tots dos.

El resultat quantitatiu principal és que \code{focus\_bgq128}, sobre 20 casos
CSMR seleccionats de manera reproduïble, redueix el bitrate mitjà de 4,0602 a
3,7013~bps respecte de \code{q204}. L'estalvi agregat és 8,84\%, mentre el PSNR
de la ROI augmenta de mitjana 0,203~dB i el fons disminueix 2,204~dB respecte
de \code{q204}, sempre amb la mateixa màscara ROI/fons. No s'ha obtingut massa
compressió; s'ha demostrat una
redistribució moderada, mesurable i controlada.

\begin{table}[H]
\centering
\footnotesize
\begin{tabularx}{\textwidth}{p{0.7cm}>{\raggedright\arraybackslash}X}
\toprule
 & Resposta final \\
\midrule
P1 & Sí, dins l'abast experimental: el pipeline C comprimeix, descomprimeix,
reconstrueix i conserva els artefactes necessaris amb memòria limitada. \\
P2 & Sí, amb dependència del context: els presets automàtics substitueixen la ROI
manual quan produeixen una cobertura interpretable, però no tots els presets són
útils en totes les escenes. \\
P3 & Sí, en la mostra CSMR: \code{focus\_bgq128} redueix bitrate respecte de
\code{q204} i conserva la ROI; \code{preserve-roi} canvia el compromís cap a més
fidelitat de ROI. \\
P4 & Parcialment: decidir el Q-map és molt barat, però el còdec domina el cost.
La ruta amb error mesurat és lenta; la taula precalibrada pot ser atractiva si la
calibració és representativa. \\
\bottomrule
\end{tabularx}
\caption{Resposta sintètica a les preguntes de recerca formulades al
Capítol~\ref{ch:introduccio}.}
\label{tab:research_question_answers}
\end{table}

\section{Compliment dels objectius}\label{sec:objective_answers}

Els cinc objectius definits al Capítol~\ref{ch:introduccio} s'avaluen de manera
explícita a la Taula~\ref{tab:objective_results}. La distinció entre
\emph{assolit} i \emph{assolit amb abast} evita presentar com a general un
resultat condicionat pel dataset o per la plataforma utilitzada.

\begin{table}[H]
\centering
\scriptsize
\begin{tabularx}{\textwidth}{p{0.7cm}p{2.1cm}X}
\toprule
 & Estat & Evidència de compliment \\
\midrule
O1 & Assolit amb abast & S'han integrat compressor, descompressor i format de
flux en C. Els tests i les comparacions amb la referència verifiquen el
comportament dins les toleràncies documentades. \\
O2 & Assolit amb abast & Els presets C identifiquen regions sense intervenció
manual i es combinen amb polítiques de qualitat. Alguns presets continuen
limitats per la cobertura del dataset i l'absència de bandes SWIR. \\
O3 & Assolit en la mostra & CSMR mesura bitrate real amb fitxers \code{.tfci}
i una màscara comuna. \code{focus\_bgq128} obté un 8,84\% d'estalvi mitjà
respecte de \code{q204} en els 20 casos principals. Aquesta xifra correspon al
codificador de referència, no al bitstream C sense codificació per entropia. \\
O4 & Assolit & L'optimització C aporta un speedup total mitjà
d'$1,205\times$ en 66 parelles Raspberry, sense fallades de paritat funcional. \\
O5 & Assolit amb abast & La decisió semàntica representa un 0,043\% del temps
de compressió, però calcular qualitat amb error mesurat depèn del còdec complet:
261,32~s per compressió i 468,13~s per compressió més descompressió en
Raspberry. L'estalvi de transmissió compensa el detector; la viabilitat de la
ruta completa depèn del throughput, de més optimització i de completar la
codificació per entropia o aritmètica en C. La taula precalibrada és una
alternativa prometedora quan la calibració és transferible, perquè evita la
passada d'error base. \\
\bottomrule
\end{tabularx}
\caption{Compliment dels objectius específics formulats al
Capítol~\ref{ch:introduccio}.}
\label{tab:objective_results}
\end{table}

\section{Avaluació de les hipòtesis}\label{sec:hypothesis_answers}

\begin{table}[H]
\centering
\scriptsize
\begin{tabularx}{\textwidth}{p{0.6cm}p{2.0cm}X}
\toprule
 & Estat & Resposta \\
\midrule
H1 & Acceptada amb abast & El pipeline C executa compressor i descompressor,
manté els formats i funciona en menys d'1~GiB. La correspondència amb la
referència s'ha validat amb operacions unitàries, artefactes i toleràncies; no
s'afirma identitat universal TensorFlow--C per qualsevol plataforma. \\
H2 & Parcialment acceptada & Diversos presets produeixen ROI interpretables
sense operador i mantenen qualitat en CSMR. No tots els presets funcionen en
totes les escenes i falta B11 per validar completament núvols. \\
H3 & Acceptada en la mostra & \code{focus\_bgq128} redueix bitrate real
respecte de \code{q204} i conserva la ROI en els casos seleccionats. El resultat no
es generalitza encara a qualsevol sensor o distribució. \\
H4 & Acceptada & \code{focus} afavoreix estalvi; \code{preserve-roi} fixa
Q240/Q255 a la ROI i afavoreix fidelitat. Cap política domina simultàniament
les dues dimensions. \\
H5 & Acceptada & Les optimitzacions de GDN/IGDN, depth-to-space i transformada
espectral aporten un speedup total mitjà d'1,205$\times$ amb zero fallades de
paritat en 66 parelles Raspberry. \\
H6 & Acceptada amb abast & Un preset semàntic tarda 0,113~s de mitjana: 0,043\%
del temps del compressor. El coll d'ampolla és el còdec, que també determina el
cost de calcular qualitat amb error mesurat. Amb una taula precalibrada, aquest
cost pot reduir-se, però només si la calibració representa la missió. \\
\bottomrule
\end{tabularx}
\caption{Resposta final a les hipòtesis formalitzades al
Capítol~\ref{ch:metodologia}.}
\label{tab:hypothesis_results}
\end{table}

\section{Contribucions assolides}\label{sec:contributions_final}

Les contribucions tècniques i metodològiques es poden resumir en:

\begin{itemize}[leftmargin=*]
  \item implementació C del pipeline de compressió i descompressió SORTENY,
  amb pesos exportats, format de bitstream documentat i validació d'operacions;
  \item recalibració coherent de la qualitat amb $\lambda_{max}=0,05$ i
  separació explícita de l'escala anterior;
  \item generadors C de Q-map uniforme, objectiu global, adaptatiu i semàntic;
  l'objectiu global s'ha validat com a mecanisme de conversió objectiu--Q amb
  una corba canònica, no com a estadística multiimatge;
  \item biblioteca de presets automàtics i llindars auditats, sense ROI
  manual en la ruta operativa;
  \item dues polítiques complementàries: \code{focus\_bgq128} i
  \code{preserve-roi};
  \item validació C-only de 120 retalls i validació de bitrate real amb CSMR
  mitjançant qualitat local per bloc;
  \item criteri de cribratge per proxies C-only, amb l'índex complet documentat
  com a eina auxiliar i no com a mètrica final;
  \item registre de procedència per a cada xifra utilitzada, indicant població,
  unitat i baseline;
  \item optimització incremental per Cortex-A53 i demostració Raspberry amb
  mesures de temps, RSS, temperatura i throttling;
  \item mesura separada del cost dels presets i del còdec, que permet una
  decisió cost--benefici quantitativa sobre calcular qualitat a bord.
\end{itemize}

La contribució no és un nou model neuronal entrenat. És convertir un model de
recerca en un sistema C controlable espacialment, verificable i executable en
una plataforma limitada.

\section{Recomanació operativa}\label{sec:operational_recommendation}

La configuració recomanada depèn de l'objectiu:

\begin{description}[leftmargin=3.5cm,style=nextline]
  \item[Fallback robust] \code{q204}. Qualitat uniforme, comportament fàcil
  d'auditar i cap dependència d'una detecció semàntica.
  \item[Mode principal] \code{preset + focus\_bgq128}. Adequat quan existeix
  una ROI interpretable i la missió vol reduir downlink mantenint-la.
  \item[ROI prioritària] \code{preset + preserve-roi Q240/Q128} o Q255/Q128.
  Adequat quan és preferible cedir part de l'estalvi per augmentar la garantia
  de qualitat a la regió seleccionada.
  \item[Control tècnic] \code{adaptive\_s8}. Útil per estudiar dificultat local,
  però sense estalvi de bitrate mesurable davant de \code{q204} en la mostra
  principal.
  \item[Baixa cadència o calibració coneguda] Taula precalibrada. Pot evitar la
  passada de compressor+descompressor per obtenir l'error base i estalviar uns
  468~s per retall, sempre que la calibració sigui representativa del sensor,
  escena i ROI esperats.
\end{description}

Aquesta recomanació assumeix que el preset selecciona una ROI interpretable. Si
la ROI és buida, total o massa dispersa, el Capítol~\ref{ch:resultats} mostra
que el valor operatiu baixa i que convé tornar a qualitat uniforme, canviar de
preset o revisar el llindar.

Els presets \code{cloud\_avoid}, \code{local\_contrast}, \code{vegetation},
\code{vegetation\_green} i \code{chlorophyll} disposen d'evidència principal.
\code{high\_ndvi} és un candidat útil després de l'auditoria experimental.
\code{low\_ndvi}, \code{dark\_regions} i \code{water\_body} necessiten més
cobertura. \code{clouds} i \code{cloud\_avoid} s'han de revalidar amb SWIR
abans d'atribuir-los precisió física de detecció de núvols.

\section{Limitacions}\label{sec:final_limitations}

La interpretació final està condicionada per les limitacions següents:

\begin{enumerate}[leftmargin=*]
  \item El dataset té 120 retalls de vuit bandes i no inclou B11/B12.
  \item Els 20 casos CSMR principals provenen d'un cribratge C-only; la mostra
  és deliberadament informativa i no una estimació aleatòria de tota la
  distribució.
  \item Els llindars s'han ajustat sobre el dataset disponible; encara falta
  separar conjunt de desenvolupament i conjunt de test per estimar-ne la
  transferibilitat.
  \item El bitstream C desa latents enters sense codificador per entropia o
  codificador aritmètic. El bitrate final es mesura amb CSMR \code{.tfci}.
  \item CSMR és la referència de bitrate de SORTENY amb TensorFlow Compression,
  però no substitueix la necessitat futura d'un codificador final completament
  en C.
  \item PSNR i MSE mesuren fidelitat, però no el rendiment d'una aplicació
  científica posterior.
  \item \code{global\_target} s'ha validat amb una corba canònica per mostrar el
  mecanisme i la saturació fora de rang; la seva precisió estadística requeriria
  una auditoria multiimatge i multiobjectiu.
  \item La Raspberry demostra execució de baix recurs, però les mesures estan
  afectades pels estats de throttling registrats i no representen una cota ideal
  del processador.
  \item La prova Raspberry no demostra temps real, consum energètic, tolerància
  a radiació ni qualificació espacial.
\end{enumerate}

Aquestes limitacions no invaliden la comparació interna perquè els controls
comparteixen retalls, màscares i còdec. Sí que limiten la generalització
externa.

\section{Treball futur prioritzat}\label{sec:future_work}

El següent increment hauria de prioritzar resultats de sistema, no afegir modes
sense una pregunta experimental clara.

\subsection{Codificació per entropia o aritmètica en C}

La prioritat és integrar un codificador per entropia o aritmètic compatible amb
els latents. Això permetria eliminar la dependència de CSMR per mesurar bitrate
real, comparar temps completament en C i estudiar el cost del codificador a
bord. El format hauria d'incloure versió, longituds, comprovació d'integritat i
límits estrictes abans de modificar la ruta operativa validada.

\subsection{Validació espectral i de dataset}

Cal repetir els presets de núvol amb B11 i ampliar estacions, cobertes i
sensors. Els llindars s'haurien d'estimar sobre un conjunt de desenvolupament
i avaluar sobre un conjunt separat. \code{water\_body} necessita específicament
un full complet a llindar 0,10; un únic crop CSMR no permet generalitzar.

\subsection{Mètriques orientades a tasca}

La ROI hauria d'avaluar-se també amb una tasca final: classificació de coberta,
segmentació, detecció d'incendis, estimació de clorofil·la o una altra variable
de missió. Una variació petita de PSNR pot ser irrellevant o crítica segons la
tasca.

\subsection{Rendiment i energia}

El throughput del Cortex-A53 és insuficient per afirmar processament en temps
real d'un sensor d'alta taxa. Cal perfilar convolucions, explorar NEON efectiu,
tiling, acceleradors o un OBC més potent, i mesurar joules per retall. La
comparació futura ha d'utilitzar el cabal del sensor i el calendari de
contactes de la missió concreta. També cal decidir formalment quan convé usar
error mesurat i quan convé usar una taula precalibrada: la decisió depèn de la
cadència, energia disponible, memòria de buffer i risc de transferir una
calibració fora del seu domini.

\subsection{Qualitat multidimensional}

Lambdes independents per bandes espectrals o grups de canals latents podrien
adaptar millor el pressupost a una tasca. Les lambdes per capes internes són una
línia més profunda: canvien la funció del model i probablement exigeixen
redisseny i reentrenament. No s'han de presentar com una extensió directa del
Q-map actual.

\section{Conclusió final}\label{sec:final_conclusion}

El treball demostra que és possible controlar espacialment un compressor
neuronal multiespectral des de C amb una interfície simple i reproduïble. La
selecció automàtica elimina la necessitat d'una ROI dibuixada per un operador,
i el cost de fer-la és negligible davant del còdec. En casos representatius,
la política focus redueix aproximadament un 8,84\% el bitrate real respecte de
qualitat uniforme, conservant la regió seleccionada i traslladant la distorsió
al fons. Preserve-roi ofereix una alternativa quan la fidelitat de la ROI té
més valor que l'estalvi màxim.

La implementació optimitzada funciona en Raspberry amb memòria limitada i és
un $1,205\times$ més ràpida en speedup total mitjà que la versió baseline, cosa
que equival a reduir el temps aproximadament un 17\%, sense canvis funcionals
detectats. La mateixa prova mostra el límit actual: la
latència del còdec continua sent de minuts per retall quan es calcula qualitat
amb error mesurat. La taula precalibrada obre una ruta més lleugera si la missió
pot assumir una calibració prèvia representativa, però no elimina la necessitat
de validar qualitat i transferibilitat. Per tant, la conclusió és positiva però
acotada: el control espacial és útil i barat; la prioritat per apropar el
sistema a una missió real és augmentar el throughput, completar la codificació
per entropia o aritmètica en C i decidir quan és acceptable substituir error
mesurat per calibració precalculada.
